\section{Дневник выполнения работы}
\TemplateNote{Журнал создан как ориентир на 24--27 ноября 2026 года.
Ниже --- плановые образцы в прошедшем времени, а не подтверждённые
записи о выполнении. Фактические действия, даты, средства и результаты
вносятся после работы. При переносе этапа изменить дату фактической записи.}

\DayRecord{24 ноября 2026 года: подготовка и начало библиотеки}
{Создал журнал выполнения лабораторной работы. Разобрал требования
к словарю и отчёту, выделил граничные случаи команд. Повторил инварианты
B-дерева, поиск и разделение узла. Определил интерфейс библиотеки,
представление записей и правила владения памятью. Начал реализацию
поиска, проверки инвариантов и вставки.}
{Условие, конспект B-дерева, редактор кода, Git;
\Field{компилятор и реальные команды}.}
{Описание интерфейса, первые исходники и тесты;
\Field{пути и коммит этапа}.}

\DayRecord{25 ноября 2026 года: операции дерева и файлы}
{Завершил вставку и обработку разделений. Реализовал заимствование,
слияние, удаление и сокращение корня. Добавил команды, нормализацию
слов и значения полного 64-битного диапазона. Определил двоичный формат,
реализовал Save и Load во временное дерево. Подготовил базовую версию
для последующего исследования качества.}
{Компилятор, ручные тесты ветвей удаления и файловых операций,
валидатор инвариантов, Git; \Field{версии и команды}.}
{Базовая версия, описание формата, тесты Save/Load;
\Field{хеш базового коммита и логи}.}

\DayRecord{26 ноября 2026 года: диагностика и исправления}
{Подготовил сравнение с эталонным контейнером и воспроизводимые случайные
тесты. Проверил граничные ключи и значения, повреждённые файлы
и сохранность словаря при ошибочном Load. Выполнил диагностику памяти
и профилирование, зафиксировал результаты. Для обнаруженных проблем
подготовил воспроизводимые примеры, внёс необходимые исправления
и повторил проверки.}
{Планируемые средства: std::map, ASan/UBSan, Valgrind Memcheck и gprof;
при необходимости Massif.
\Field{фактически использованные средства, версии и команды}.}
{Тесты, seed и исходные логи, описание реальных проблем или объёма
проверок без найденных проблем, исправленная версия;
\Field{коммиты и пути к подтверждениям}.}

\DayRecord{27 ноября 2026 года: измерения и оформление}
{Подготовил одинаковые нагрузки для сравниваемых версий. Зафиксировал среду
и параметры сборки. Выполнил повторные измерения времени и памяти,
сопоставил результаты, оформил таблицы и объяснил различия.
Уточнил дневник по фактически выполненной работе, сформулировал
выводы и собрал итоговый отчёт.}
{Обычная оптимизированная сборка, средство замера времени и пиковой памяти,
отдельные диагностические сборки, LaTeX и Git;
\Field{команды, версии и число повторов}.}
{Таблицы и сырые измерения, сравниваемые коммиты, заполненный отчёт;
\Field{логи, PDF и финальная версия}.}
